\originalchapter{平面图与对偶}\label{ch:planar}
\sourceof{6.042J 第12章; 18.315 L6--7、L10; 桥与面边界计数补全}
\originalsection{嵌入与面}
\begin{definition}
平面图是一张已经在平面中画出的图, 边的内部彼此不交且不经过其他顶点. 可平面图是存在这种画法的抽象图. 平面去掉所画顶点与边后的连通区域称为面, 包括唯一无界外面. 面的度为沿其全部边界行走得到的边出现总次数.
\end{definition}
同一抽象图可有不同嵌入, 但连通图的面数由顶点与边数决定. 桥两侧属于同一个面, 在面边界中计两次. 因而即使一个面不是由简单圈围成, 仍有面度总和 $\sum_F d(F)=2m$.

以下采用平面拓扑中的 Jordan 圈分离定理作为外部先修工具: 一条简单闭曲线把球面分成两个连通区域. 本书不证明该拓扑定理, 而是明确用它支持“圈边两侧为不同面”以及五色证明中的不能交叉. 平面加上无穷远点即为球面, 外面与其他面于是平等.
\begin{theorem}[Euler 公式]\label{thm:eulerformula}
平面嵌入图有 $n$ 个顶点、$m$ 条边、$f$ 个面和 $c$ 个分支时, 满足 $n-m+f=1+c$.
\end{theorem}
\begin{proof}
连通图若无圈, 是树, 有 $m=n-1$, 且补集只有一个面, 公式成立. 若有圈, 删除圈上的一边. 两侧面不同, 删除边将它们合并, 所以边数和面数各减1; 图仍连通. 反复删除直到树, $n-m+f$ 始终不变, 因而为2.

非连通图可在同一个面的不同分支之间加边, 连接 $c-1$ 次后连通. 新边是桥, 不分割面, 面数不变. 对新图用连通公式: $n-(m+c-1)+f=2$, 整理得到结论.
\end{proof}
\begin{corollary}\label{cor:planaredges}
简单可平面图若 $n\ge3$, 则 $m\le3n-6$. 若还无三角形, 则 $m\le2n-4$. 因而 $K_5$ 和 $K_{3,3}$ 都不可平面.
\end{corollary}
\begin{proof}
先设连通. 当 $n\ge3$ 且简单时, 每个面边界长度至少3; 长度1需要自环, 长度2需要平行边或整图只是一条边. 所以 $3f\le2m$, 联合 $f=2-n+m$ 得第一界. 无三角形时每面边界长度至少4, 得第二界. 树的边界可能重复边, 但在 $n\ge3$ 时长度 $2(n-1)\ge4$, 与计算一致.

非连通图在面内加桥使连通, 不产生圈, 因而也不产生三角形, 再应用两界. $K_5$ 有10边而 $3n-6=9$; $K_{3,3}$ 二部而无三角形, 有9边而 $2n-4=8$.
\end{proof}
少于界限也不保证可平面. 例如把 $K_{3,3}$ 每条边细分一次, 边数与顶点数增加后仍不可平面, 但可以满足 $m\le3n-6$. 细分不改变是否能无交叉画图.
\originalsection{五色定理}
\begin{theorem}[五色定理]\label{thm:fivecolor}
每个有限简单可平面图都可正常地用五色染色.
\end{theorem}
\begin{proof}
对顶点数归纳. 平面边数界与握手引理保证至少一点 $v$ 度至多5. 删除它后归纳染色. 若邻点不足5个或所用颜色不足5种, 给 $v$ 选未用色即可. 否则按嵌入中围绕 $v$ 的循环顺序, 五邻点为 $v_1,\ldots,v_5$, 其颜色依次为1至5.

考虑仅有1、3两色的诱导子图. 若 $v_1,v_3$ 不在同一分支, 交换含 $v_1$ 分支的1、3两色, 保持正常染色并在 $v$ 邻域释放颜色1. 若它们相连, 取一条简单双色路, 加上 $vv_1,vv_3$ 形成简单闭曲线. 邻点 $v_2,v_4$ 位于其不同侧. 2、4双色路不能经过1、3双色路的顶点, 也不能穿越其边, 因此 $v_2,v_4$ 在2、4诱导子图中不连通. 交换含 $v_2$ 分支的2、4两色, 释放颜色2, 再给 $v$ 染色.
\end{proof}
四色定理进一步把五色改为四色, 但其证明不属于本册. 五色定理不是四色定理的证明, 两条双色链能否同时交换必须按平面分离关系检查.
\begin{center}
\begin{tikzpicture}[scale=1.05]
\node[vertex](v)at(0,0){$v$};
\foreach \i/\ang in {1/90,2/18,3/-54,4/-126,5/162}{\node[vertex](v\i)at(\ang:1.8){$v_{\i}$};\draw[edge](v)--(v\i);}
\draw[very thick,color=structurecolor](v1)..controls(3.4,2.6)and(3.4,-2.3)..node[pos=.48,right=3mm]{$1,3$路}(v3);
\node at(0,-2.5){循环顺序决定 $v_2,v_4$ 分处两侧};
\end{tikzpicture}
\end{center}
\originalsection{平面对偶}
\begin{definition}
连通平面图 $G$ 的对偶图 $G^*$ 在每个面中放一个顶点, 对每条原边加一条穿过它一次的对偶边, 连接其两侧的面. 对偶通常是允许自环和平行边的多重图. 原边与对偶边一一对应.
\end{definition}
\begin{proposition}\label{prop:dual-tree}
若 $T$ 为 $G$ 的生成树, 则原图不在 $T$ 中的边对应的对偶边组成 $G^*$ 的生成树. 原图的桥对应对偶自环.
\end{proposition}
\begin{proof}
桥两侧同一面, 因而对应自环. 对生成树, 先只画树, 补集连通. 逐条加入原图的非树边; 每次新边都与现有图中的路构成圈, 把一个面分成两个. 在追踪面分裂的对偶记录中, 新边连接这两个子面, 相当于把一个已有顶点分裂成两个并以新边连接. 从单点开始这样构造, 所得记录始终是一棵树. 最终它恰为全部非树边的对偶边, 连接全部 $f$ 个面, 边数为 $m-(n-1)=f-1$.
\end{proof}
\originalsection{曲面上的上界}
平面是球面去掉一点. 若允许把图画在有 $g$ 个把手的闭可定向曲面上, 圈分割关系就不同. 只在这里补充原课中的上界, 不把完整曲面分类或 Heawood 最优性证明纳入核心.
\begin{proposition}
若简单图可嵌入亏格 $g\ge1$ 的闭可定向曲面, 则
\[
\chi(G)\le H(g):=\left\lfloor\frac{7+\sqrt{1+48g}}2\right\rfloor.
\]
\end{proposition}
\begin{proof}
所用曲面 Euler 工具为 $n-m+f=2-2h$, 适用于每个面为圆盘的连通嵌入. 一般嵌入可取图的窄带邻域, 给每条边界圈封圆盘, 得亏格 $h\le g$ 的这种嵌入. 面边界长度至少3时得到 $m\le3n-6+6g$; 对 $n\ge3$ 的简单连通图仍适用, 与平面计数同理.

若染色数 $k\le7$, 因 $H(g)\ge7$ 已有结论. 若 $k\ge8$, 取第\ref{ch:coloring}章习题中的 $k$ 临界子图, 它有 $n\ge k$ 点且最小度至少 $k-1$, 并且连通: 否则某个较小分支已具有染色数 $k$, 与临界选择矛盾. 嵌入限制在子图上仍可用上述边界, 故 $(k-1)n\le6n-12+12g$, 即 $(k-7)n\le12g-12$. 因 $k-7>0$, 有 $k(k-7)\le12g-12$, 解二次式得 $k\le(7+\sqrt{1+48g})/2$, 再取整.
\end{proof}
曲面 Euler 公式和“封边界不增加亏格”在这里作为明确的拓扑工具引用. 球面 $g=0$ 的四色结论不由这个推导给出. Kuratowski 的两个细分禁图判据以及 Tutte 的四连通平面图 Hamilton 定理也只列为进一步阅读, 不作为正文后续推理的前提.
\originalsection{习题与解答}
\begin{exercise}\label{ex:plan1}
若连通简单平面图的每个面都为三角形, 求边数与面数.
\end{exercise}
\begin{exercise}\label{ex:plan2}
证明每个简单二部可平面图都有度数至多3的顶点, 并由此给出四色染色算法.
\end{exercise}
习题\ref{ex:plan1}的解答.
\begin{solution}
面度和为 $3f=2m$, Euler 给 $n-m+f=2$, 得 $m=3n-6$, $f=2n-4$. 条件包括外面也为三角形, 不可漏算它.
\end{solution}
习题\ref{ex:plan2}的解答.
\begin{solution}
$n\ge3$ 时无三角形边界给平均度 $2m/n\le4-8/n<4$, 所以有度至多3点; 小图直接成立. 每个诱导子图仍二部且可平面, 可不断删度至多3点, 再逆序贪心用四色. 实际由二部性只需两色; 此处算法用于说明边数界如何给退化上界.
\end{solution}
